% ============================================================
% appendix_interview.tex — Appendix: Interview Questions
% ============================================================

\chapter{Questions Every High Performer Asks in an Interview}

\sectionintro{``An interview goes both ways. You're not desperately
trying to convince them to hire you. You're trying to decide whether
you should join.''}

The best candidates presume they will get the job. They ask questions
that demonstrate they aren't going to accept just any offer. This
counter-intuitively makes the interviewer want to hire them more.

The interview isn't over when they ask, ``Do you have any questions
for me?''

\vspace{1em}

\begin{enumerate}[leftmargin=2em, itemsep=1.2em]

  \item \textbf{If I were hired, what's the first thing I would work on?}

  Helps you understand what the team's actual needs are, and plants a
  seed by having the interviewer imagine you already on the team.

  \item \textbf{I noticed your team just launched \textit{X}.
  What have you learned from the launch?}

  Demonstrates that you've done your homework. Gives you insight into
  how people think about their own work --- and whether they reflect
  on it.

  \item \textbf{What are the KPIs for the team, org, or company?
  Have you hit your targets in the last cycle?}

  Gives you a sense of how organizational performance is measured ---
  and whether the team is honest about missing it.

  \item \textbf{How much attrition has the team experienced this year,
  and what are the primary drivers?}

  Most candidates ask about growth. Asking about departures paints a
  more accurate picture of the team's reality.

  \item \textbf{How many reorgs have you experienced in the last year?}

  Gives you a sense of how much organizational churn the team has
  absorbed --- and what it costs them in focus.

  \item \textbf{How much of the roadmap have you shipped in the last
  six months?}

  Helps you understand whether the team you're about to join is in
  maintenance mode or growth mode.

\end{enumerate}

\vspace{1.5em}

The best candidates presume they will get the job. The questions they
ask prove it.
